\scanpage{295}
- \mspage{11} -
\margin{$n^\circ\ 246$}

similitude amenant $D$ sur $D'$ s'appellera le \uline{diangle} de $D$ et $D'$ et se notera $(D,D')$. Le diangle déterminé par la similitude $S$ est d'ordre $2$ dans le groupe $G_{\unsure{0|o}}/H$. Pour qu'un élément $X=\xi\cdot1_E+\eta S$ de $G_0$ détermine un diangle d'ordre $2$, il faut et il suffit qu'on ait $2\xi\eta=0$ ; donc le diangle déterminé par $S$ est le seul élément d'ordre $2$ dans $G_0/H$ si $k$ n'est pas de caractéristique $2$.

Si $X=\xi\cdot1_E+\eta S$ est un élément de $G_0$, $\xi,\eta$ ne sont pas tous deux nuls; posons $\tau(X)=\eta/\xi$ si $\xi\ne0$, et $\tau(X)=\infty$ si $\xi=0$ ; autrement dit, on pose en tout cas $\tau(X)=\eta/\xi$, le second membre étant entendu au sens du \uline{corps projectif} $k_\infty$ attaché à $k$ (Alg., Chap. II, 2e éd., App. III, $n^\circ\ 5$); Pour que deux éléments $X,X'$ de $G_0$ appartiennent à la même classe suivant $H$, il faut et il suffit que $\tau(X)=\tau(X')$ ; pour qu'un élément $\tau$ de $k_\infty$ soit image par $\tau(X)$ d'un élément de $G_0$, il faut et il suffit qu'on ait $\delta\tau^2\ne1$. Autrement dit, $\tau(X)$ détermine une \uline{bijection} de $G_0/H$ sur le complémentaire dans $k_\infty$ de l'ensemble des racines de $\delta\tau^2=1$ (ensemble à deux éléments si $\delta$ est un carré \unsure{d'ns|dans} $k$, vide dans le cas contraire). On a la formule :
\[
\tau(XX')=\frac{\tau(X)+\tau(X')}{1+\delta\tau(X)\tau(X')}\qquad(9),
\]
dont la vérification est immédiate, d'abord dans le cas où $\tau(X),\tau(X')$, $\tau(XX')$ sont finis, puis dans tous les autres cas.

L'homomorphisme canonique de $G_0$ sur $G_0/H$ induit sur $\Gamma_{\unsure{0|o}}$ un homomorphisme de $\Gamma_0$ dans $G_0/H$ dont le noyau est $\Gamma_0\cap H=\{\pm1_E\}$, groupe à deux éléments si $k$ n'est pas de caractéristique $2$. Mais dans ce cas on peut aussi définir un isomorphisme (non moins \unsure{khânonique|khanonique}) entre ces groupes :

PROPOSITION 2. -- \uline{Pour tout} $X\in G_0$, $N(X)^{-1}X^2$ \uline{est une rotation} ;
\mspage{294}
